% GENERATED FILE. Edit canonical lessons, then run npm run generate:books.
% canonical-source-hash: fnv1a64:a800e6de19cda20f
% canonical-lessons: MW-C173-phukno, MW-C173-dhuvo, MW-C173-idhan, MW-C173-bujhavno, MW-C173-sekno

\chapter{At the Hearth}
\label{ch:mw-at-the-hearth}
\hlchaptermodality{173}

\begin{chapteropening}
\textbf{By the end of this chapter:} \emph{I can say how the fire is lit, used and put out.}

Complete the last lesson of chapter 173: I can say how the fire is lit, used and put out.
\end{chapteropening}

\section[\texorpdfstring{\mw{फूंकणो} (phū̃kṇo)}{फूंकणो (phū̃kṇo)}]{\mw{फूंकणो} (phū̃kṇo) — to blow on (a fire)}

\label{lesson:MW-C173-phukno}



\begin{quote}
Before the new one: say the Marwadi for to plaster, then the Marwadi for mud plaster.
\end{quote}

\subsection*{You'll want to know: \mw{फूंकणो}}
\textbf{\mw{फूंकणो}} — \emph{phū̃kṇo} — \textquotedblleft{}to blow on (a fire)\textquotedblright{}.

\begin{quote}
\mw{दादी} \mw{आग} \mw{फूंकै} \mw{है।}
\end{quote}

\begin{quote}
\emph{dādī āg phū̃kai hai.}
\end{quote}

\begin{quote}
Grandmother blows on the fire.
\end{quote}

Past (I …): \mw{म्हैं} \mw{आग} \mw{फूंकी}, \emph{mhaĩ āg phū̃kī}.

\begin{cousinweb}
\textbf{\mw{फूंकणो}} (\emph{phū̃kṇo}) is inherited from a reconstructed \emph{*phūtka}, \textquotedblleft{}blowing\textquotedblright{}, a word that imitates the sound of a puff; CDIAL lists Old Marwari \emph{phūkaï}, \textquotedblleft{}to blow\textquotedblright{} (CDIAL 9102).
\end{cousinweb}

\begin{grammarlens}[title={What you've built}]
One more word for saying what you do.
\end{grammarlens}

\subsection*{Your turn}
\begin{itemize}
\raggedright
  \item \emph{Say it:} \emph{phū̃kṇo}
  \item \emph{Say it:} \emph{phū̃kṇo}, once more
  \item \emph{Say it:} \emph{dādī āg phū̃kai hai.} — Grandmother blows on the fire.
  \item \emph{Recall:} say the Marwadi for to plaster, then the Marwadi for mud plaster, then say \emph{phū̃kṇo} again
\end{itemize}

\begin{tcolorbox}[breakable,colback=teal!4,colframe=teal!35!black,title={Before you move on}]
What does \mw{फूंकणो} mean? (\textquotedblleft{}To blow on (a fire)\textquotedblright{}.) Say it once more.
\end{tcolorbox}
\section[\texorpdfstring{\mw{धूंवो} (dhū̃vo)}{धूंवो (dhū̃vo)}]{\mw{धूंवो} (dhū̃vo) — smoke}

\label{lesson:MW-C173-dhuvo}



\begin{quote}
Before the new one: say the Marwadi for mud plaster, then the Marwadi for to blow on.
\end{quote}

\subsection*{You'll want to know: \mw{धूंवो}}
\textbf{\mw{धूंवो}} — \emph{dhū̃vo} — \textquotedblleft{}smoke\textquotedblright{}.

\begin{quote}
\mw{रसोई} \mw{मांय} \mw{घणो} \mw{धूंवो} \mw{है।}
\end{quote}

\begin{quote}
\emph{rasoī mā̃y ghaṇo dhū̃vo hai.}
\end{quote}

\begin{quote}
There's a lot of smoke in the kitchen.
\end{quote}

A masculine noun.

\begin{cousinweb}
\textbf{\mw{धूंवो}} (\emph{dhū̃vo}) is inherited from Sanskrit \emph{dhūma}, \textquotedblleft{}smoke, vapour\textquotedblright{}; CDIAL lists Hindi \emph{dhūā̃} (CDIAL 6849). CDIAL gives no Marwari form in this entry.
\end{cousinweb}

\begin{grammarlens}[title={What you've built}]
One more word for saying what you do.
\end{grammarlens}

\subsection*{Your turn}
\begin{itemize}
\raggedright
  \item \emph{Say it:} \emph{dhū̃vo}
  \item \emph{Say it:} \emph{dhū̃vo}, once more
  \item \emph{Say it:} \emph{rasoī mā̃y ghaṇo dhū̃vo hai.} — There's a lot of smoke in the kitchen.
  \item \emph{Recall:} say the Marwadi for mud plaster, then the Marwadi for to blow on, then say \emph{dhū̃vo} again
\end{itemize}

\begin{tcolorbox}[breakable,colback=teal!4,colframe=teal!35!black,title={Before you move on}]
What does \mw{धूंवो} mean? (\textquotedblleft{}Smoke\textquotedblright{}.) Say it once more.
\end{tcolorbox}
\section[\texorpdfstring{\mw{ईंधण} (ī̃dhaṇ)}{ईंधण (ī̃dhaṇ)}]{\mw{ईंधण} (ī̃dhaṇ) — firewood; fuel}

\label{lesson:MW-C173-idhan}



\begin{quote}
Before the new one: say the Marwadi for to blow on, then the Marwadi for smoke.
\end{quote}

\subsection*{You'll want to know: \mw{ईंधण}}
\textbf{\mw{ईंधण}} — \emph{ī̃dhaṇ} — \textquotedblleft{}firewood; fuel\textquotedblright{}.

\begin{quote}
\mw{ईंधण} \mw{सूको} \mw{है।}
\end{quote}

\begin{quote}
\emph{ī̃dhaṇ sūko hai.}
\end{quote}

\begin{quote}
The firewood is dry.
\end{quote}

A masculine noun.

\begin{cousinweb}
\textbf{\mw{ईंधण}} (\emph{ī̃dhaṇ}) is inherited from Sanskrit \emph{indhana}, \textquotedblleft{}kindling, fuel\textquotedblright{}; CDIAL lists Hindi \emph{ī̃dhan} and Gujarati \emph{ĩdhaṇ} (CDIAL 1584). CDIAL gives no Marwari form in this entry.
\end{cousinweb}

\begin{grammarlens}[title={What you've built}]
One more word for saying what you do.
\end{grammarlens}

\subsection*{Your turn}
\begin{itemize}
\raggedright
  \item \emph{Say it:} \emph{ī̃dhaṇ}
  \item \emph{Say it:} \emph{ī̃dhaṇ}, once more
  \item \emph{Say it:} \emph{ī̃dhaṇ sūko hai.} — The firewood is dry.
  \item \emph{Recall:} say the Marwadi for to blow on, then the Marwadi for smoke, then say \emph{ī̃dhaṇ} again
\end{itemize}

\begin{tcolorbox}[breakable,colback=teal!4,colframe=teal!35!black,title={Before you move on}]
What does \mw{ईंधण} mean? (\textquotedblleft{}Firewood; fuel\textquotedblright{}.) Say it once more.
\end{tcolorbox}
\section[\texorpdfstring{\mw{बुझावणो} (bujhāvṇo)}{बुझावणो (bujhāvṇo)}]{\mw{बुझावणो} (bujhāvṇo) — to put out (a fire or a lamp)}

\label{lesson:MW-C173-bujhavno}



\begin{quote}
Before the new one: say the Marwadi for smoke, then the Marwadi for firewood.
\end{quote}

\subsection*{You'll want to know: \mw{बुझावणो}}
\textbf{\mw{बुझावणो}} — \emph{bujhāvṇo} — \textquotedblleft{}to put out (a fire or a lamp)\textquotedblright{}.

\begin{quote}
\mw{रात} \mw{मांय} \mw{हूं} \mw{दीयो} \mw{बुझाऊं} \mw{हूं।}
\end{quote}

\begin{quote}
\emph{rāt mā̃y hū̃ dīyo bujhāū̃ hū̃.}
\end{quote}

\begin{quote}
At night I put out the lamp.
\end{quote}

Past (I …): \mw{म्हैं} \mw{दीयो} \mw{बुझायो}, \emph{mhaĩ dīyo bujhāyo}.

\begin{cousinweb}
\textbf{\mw{बुझावणो}} (\emph{bujhāvṇo}) is built on a reconstructed \emph{*vijjhāyati}, \textquotedblleft{}is burnt out, goes out\textquotedblright{}; CDIAL lists Old Marwari \emph{bujhāvaï}, \textquotedblleft{}to put out\textquotedblright{} (CDIAL 11703).
\end{cousinweb}

\begin{grammarlens}[title={What you've built}]
One more word for saying what you do.
\end{grammarlens}

\subsection*{Your turn}
\begin{itemize}
\raggedright
  \item \emph{Say it:} \emph{bujhāvṇo}
  \item \emph{Say it:} \emph{bujhāvṇo}, once more
  \item \emph{Say it:} \emph{rāt mā̃y hū̃ dīyo bujhāū̃ hū̃.} — At night I put out the lamp.
  \item \emph{Recall:} say the Marwadi for smoke, then the Marwadi for firewood, then say \emph{bujhāvṇo} again
\end{itemize}

\begin{tcolorbox}[breakable,colback=teal!4,colframe=teal!35!black,title={Before you move on}]
What does \mw{बुझावणो} mean? (\textquotedblleft{}To put out (a fire or a lamp)\textquotedblright{}.) Say it once more.
\end{tcolorbox}
\section[\texorpdfstring{\mw{सेकणो} (sekṇo)}{सेकणो (sekṇo)}]{\mw{सेकणो} (sekṇo) — to toast (on a fire or a griddle)}

\label{lesson:MW-C173-sekno}



\begin{quote}
Before the new one: say the Marwadi for firewood, then the Marwadi for to put out.
\end{quote}

\subsection*{You'll want to know: \mw{सेकणो}}
\textbf{\mw{सेकणो}} — \emph{sekṇo} — \textquotedblleft{}to toast (on a fire or a griddle)\textquotedblright{}.

\begin{quote}
\mw{मां} \mw{आग} \mw{माथै} \mw{पापड़} \mw{सेकै} \mw{है।}
\end{quote}

\begin{quote}
\emph{mā̃ āg māthai pāpaṛ sekai hai.}
\end{quote}

\begin{quote}
Mother toasts a papad over the fire.
\end{quote}

Past (I …): \mw{म्हैं} \mw{रोटी} \mw{सेकी}, \emph{mhaĩ roṭī sekī}.

\begin{cousinweb}
\textbf{\mw{सेकणो}} (\emph{sekṇo}) is inherited from a reconstructed \emph{*sekk}, \textquotedblleft{}heat, roast\textquotedblright{}; CDIAL lists Old Marwari \emph{sekaï}, \textquotedblleft{}heats, roasts\textquotedblright{} (CDIAL 13581).
\end{cousinweb}

\begin{grammarlens}[title={What you've built}]
One more word for saying what you do.
\end{grammarlens}

\subsection*{Your turn}
\begin{itemize}
\raggedright
  \item \emph{Say it:} \emph{sekṇo}
  \item \emph{Say it:} \emph{sekṇo}, once more
  \item \emph{Say it:} \emph{mā̃ āg māthai pāpaṛ sekai hai.} — Mother toasts a papad over the fire.
  \item \emph{Recall:} say the Marwadi for firewood, then the Marwadi for to put out, then say \emph{sekṇo} again
\end{itemize}

\begin{tcolorbox}[breakable,colback=teal!4,colframe=teal!35!black,title={Before you move on}]
What does \mw{सेकणो} mean? (\textquotedblleft{}To toast (on a fire or a griddle)\textquotedblright{}.) Say it once more.
\end{tcolorbox}
